\documentclass[10pt,oneside]{book}
%%%% Page Info + Commands %%%%%{

%packages
\usepackage{pgfplots}
\usepackage{mathrsfs}
\pgfplotsset{compat=1.18}
\usepackage{amsmath}
\usepackage{amsfonts}
\usepackage{charter}
\usepackage{amsthm,letltxmacro}
\usepackage{amssymb}
\usepackage{fancyhdr}
\usepackage{float}
\usepackage[most]{tcolorbox}
\usepackage{enumerate}
\usepackage{enumitem}
\usepackage[dvipsnames]{xcolor}
\usepackage{changepage}
\usepackage{graphicx}

% SMALL PAGE COMMAND
\newcommand{\smallpage}{  \setlength \oddsidemargin{.25in}
            \setlength \textwidth{5.5in}
            \setlength \topmargin{0in}
            \setlength \textheight{9in}}


% Commands for sets of numbers
\newcommand{\Z}{\mathbb{Z}}\newcommand{\R}{\mathbb{R}}\newcommand{\N}{\mathbb{N}}\newcommand{\Q}{\mathbb{Q}}\newcommand{\C}{\mathbb{C}}
\newcommand{\real}[1]{\text{Re}\left(#1\right)}
\newcommand{\im}[1]{\text{Im}\left(#1\right)}

% gordie spacing and boxing commands
\newcommand{\gspc}{\vspace{0.13cm}} % 0.5 space
\newcommand{\gline}[0]{\gspc\\} % 1.5 space
\newcommand{\gbline}{\gspc\gspc\\} % double space
\newcommand{\gbox}[1]{\fbox{\parbox{\linewidth}{#1}}} % multiline box command
\newcommand{\gmod}[1]{\,(\text{mod}\,#1)}


\newcommand{\gimage}[2]{
\begin{figure}[H]
    \centering
    \includegraphics[width=#2]{#1}
\end{figure}
}

\newcommand{\centerit}[1]{{\begin{center} #1 \end{center}}}
\newcommand{\ul}[1]{\underline{#1}}

\newcommand{\prt}[2]{\frac{\partial#1}{\partial#2}}

%%%%%%%% command for graphics %%%%%%%%%%%%%

%}

% COLOR BOX

\tcbset{problembox/.style={
    colback=gray!4,
    colframe=gray!50, 
	  boxrule=0pt,
    leftrule=2.5pt, 
    left=2mm, right=2mm, top=1mm, bottom=1mm,
    boxsep=2mm,
    breakable,
    sharp corners
  }
}

\tcbset{subproblembox/.style={
    colback=gray!12,
    colframe=gray!80, 
	  boxrule=0.5pt,
    leftrule=2.5pt, 
    left=2mm, right=2mm, top=1mm, bottom=1mm,
    boxsep=2mm,
    breakable,
    sharp corners
  }
}

% PROBLEM
\newenvironment{problem}[1]
  {\vspace{-0.1cm}\begin{tcolorbox}[problembox]
   \textit{Proof. }#1}
  {\end{tcolorbox}\gspc}

  \newenvironment{subproblem}[2]
  {\vspace{-0.1cm}\begin{tcolorbox}[subproblembox]
   \textit{#1 }#2}
  {\end{tcolorbox}\gspc}





% color commands
\newcommand{\cg}[1]{\color{OliveGreen}#1\color{white}}
\newcommand{\cb}[1]{\color{BlueViolet}#1\color{white}}


\newcommand{\cpr}[1]{\left(#1\right)}
%%%%%%%%%%%%% PHYSICS SPECIFIC COMMANDS

\newcommand{\vA}{\vec{\mathbf{A}}}
\newcommand{\vB}{\vec{\mathbf{B}}}
\newcommand{\vC}{\vec{\mathbf{C}}}
\newcommand{\vZ}{\vec{\mathbf{0}}}
\newcommand{\norm}{\vec{\mathbf{n}}}
\newcommand{\pvec}[1]{\left\langle#1\right\rangle}
\newcommand{\ar}{\vec{\mathbf{r}}}
\newcommand{\arhat}{\hat{\mathbf{r}}}
\newcommand{\vv}{\vec{\mathbf{v}}}
\newcommand{\delfull}{\pvec{\prt{}{x}, \prt{}{y},\prt{}{z}}}
\newcommand{\delinput}[3]{\pvec{\prt{}{x}\cpr{#1}, \prt{}{y}\cpr{#2},\prt{}{z}\cpr{#3}}}
\newcommand{\crossprod}[6]{\text{Det}\left|\begin{matrix}\hat{\textbf{i}}&\hat{\textbf{j}} &\hat{\textbf{k}}\\#1 & #2 & #3\\#4 & #5&#6\end{matrix}\right|}
\newcommand{\dps}{\displaystyle}
\newcommand{\delmatrix}[3]{\begin{bmatrix}\dps\hat x&\dps\hat y&\dps\hat z\\\dps\prt{}x&\dps\prt{}y&\dps\prt{}z\\\dps#1&\dps#2&\dps#3\end{bmatrix}}

%%%%%%%%%%%%%%%%%%%%%%%%%%%%
%%%%%%%%%%%%%%%%%%%%%%%%%%%%%%
%%%%%%%%%%%%%%%%%%%%%%%%%%%%%
%%%%%%%%%%%%%%%%%%%%%%%%%%%%%
%%%%%%%%%%%%%%%%%%%%%%%%%%%%%%
%%%%%%%%% begin doc

%%%% Page Setup %%%%%%%
\smallpage\sloppy
\pagestyle{fancy}
\lhead{\large{\textbf{HW5 - Theory of Computation}}}
\setlength{\headheight}{10pt}


\begin{document}

\subsection*{Problem 1}
Let $A$ be the language of strings with twice as many 1's as 0's over the alphabet $\{0, 1\}$. Use the pumping lemma for regular langauges to prove that $A$ is not regular. 
\begin{problem}
  Assume for the sake of contradiction that $A$ is a regular language with pumping length $p$. Thus, for a string $s = 1^{2p}0^p\in A$, we must have an $xyz = s$ satisfying:
  \begin{itemize}[itemsep=0pt, parsep=0pt]
    \item $|y|> 0$, 
    \item $|xy| \le p$, and
    \item $xy^iz \in A$ for $i \in \N$. 
  \end{itemize}
  However, because $|y|>0$ and $|xy| < p$, we must have that $x$ and $y$ are made exclusively of 1's. Thus, although $xyz \in A$, $xy^0z\notin A$ because there are fewer than two times the number of 1s than 0s. \gline{}
  Hence, $A$ is not regular.\qed
\end{problem}

\subsection*{Problem 2}
Let $B$ be the language of strings with more 1's than 0's over the alphabet $\{0, 1\}$. Prove that $B$ is not regular. 
\begin{problem}
  Assume for the sake of contradiction that $B$ is a regular language with pumping length $p$. Consider the string $\displaystyle s = 0^p\,1\,1^p\in B$. Because $B$ is regular, there must exist an $xyz = s$ satisfying:
  \begin{itemize}[itemsep=0pt, parsep=0pt]
    \item $|y|> 0$, 
    \item $|xy| \le p$, and
    \item $xy^iz \in A$ for $i \in \N$. 
  \end{itemize}
  However, because $|xy|\le p$, we know that $x$ and $y$ are made up of exclusively 0's. Furthermore, because $|y|\ge 1$, we know $y$ must contain at least one 0. Thus, the string:
  \begin{align*}
    xy^2z = 0^{p+n}1^{p+1}\notin A, \text{ with }n \ge 1
  \end{align*}
  The number of ones is not greater than the number of zeros, and thus the string does not satisfy the pumping lemma and $B$ is not regular.  
\end{problem}

\subsection*{Problem 3}
Draw the following pushdown automata for each language:
\begin{enumerate}
  \item [a.] $\{0^k1^{2k}\mid k\ge 0\}$ This language consists of strings with zero or more 0's, followed by twice as many 1's.
  \gimage{image.png}{12cm} 
  \item [b.] $\{0^k1^l\mid l \ge k\}$
  \gimage{image2.png}{12cm}
\end{enumerate}

\subsection*{Problem 4}
Consider the following pushdown automata with diagrams below. For each diagram describe the language that is recognized by the pushdown automata. 
\begin{enumerate}
  \item [a.] The language that is described by a. is  $\{(00) | (11)\}^*$. Essentially, its just where all 0s come in strings of consecutive even counts, and all 1s come in strings of consecutive even counts e.g (0000), (111111), but not (000). 
  \item [b.] The language that is described by b. is the language that contains an even number of characters. 
\end{enumerate}

\end{document}